\documentclass[11pt,a4paper]{article}
\usepackage[T1]{fontenc}
\usepackage[utf8]{inputenc}
\usepackage[english]{babel}
\usepackage{lmodern}
\usepackage{microtype}
\usepackage{geometry}
\geometry{margin=2.5cm}
\usepackage{hyperref}
\hypersetup{
    colorlinks=true,
    linkcolor=blue!70!black,
    urlcolor=blue!70!black,
    citecolor=blue!70!black
}
\usepackage{booktabs}
\usepackage{longtable}
\usepackage{enumitem}
\usepackage{xcolor}
\usepackage{titlesec}
\usepackage{fancyhdr}
\usepackage{graphicx}
\usepackage{listings}

\definecolor{codegray}{rgb}{0.95,0.95,0.95}
\definecolor{doctum}{rgb}{0.10,0.30,0.60}

\lstset{
  backgroundcolor=\color{codegray},
  basicstyle=\small\ttfamily,
  breaklines=true,
  frame=single,
  rulecolor=\color{gray!40},
  xleftmargin=1em,
  framexleftmargin=0.5em
}

\pagestyle{fancy}
\fancyhf{}
\fancyhead[L]{\textcolor{doctum}{\small\textbf{Doctum Consilium}}}
\fancyhead[R]{\small\nouppercase{\leftmark}}
\fancyfoot[C]{\thepage}

\titleformat{\section}{\large\bfseries\color{doctum}}{}{0em}{}[\titlerule]
\titleformat{\subsection}{\normalsize\bfseries\color{doctum!80}}{}{0em}{}


\begin{document}

\title{
  \textcolor{doctum}{\Large\textbf{Market Screener}} \\[0.5em]
  \large Multi-Source Stock Screener with Airflow Pipeline and AI Analysis \\[0.3em]
  \normalsize Technical Foundations and Architecture
}
\author{Doctum Consilium \\ \small\textit{Generated 2026-05-08}}
\date{}
\maketitle
\thispagestyle{fancy}

\begin{abstract}
Automated stock screener with intraday and nightly Airflow pipelines, FastAPI backend, Yahoo Finance and Alpaca connectors, and AI-powered analysis.
This document covers the theoretical foundations, architectural decisions,
and key technical concepts required to understand, contribute to, and
maintain this repository.
\end{abstract}

\tableofcontents
\newpage

\section{Overview}

Market Screener automates the identification of trading opportunities by
running configurable screening rules across a universe of instruments.
Two Airflow pipelines run on distinct schedules: intraday (real-time signals)
and nightly (batch scoring and database refresh).

The backend exposes a REST API for queried results, and the AI analysis
module generates natural-language summaries of screener output.


\section{Architecture}

\begin{verbatim}
Airflow Scheduler
├── Intraday DAG  → Yahoo Finance / Alpaca fetch
│                 → Apply screening rules
│                 → Write to PostgreSQL
└── Nightly DAG   → Historical scoring
                  → AI analysis (LiteLLM)
                  → Dashboard refresh

FastAPI backend ──► PostgreSQL ──► React frontend
\end{verbatim}


\section{Key Technical Concepts}
\begin{itemize}[leftmargin=1.5em]
  \item Technical indicators: RSI $= 100 - \frac{100}{1 + RS}$ where $RS = \frac{\text{avg gain}}{\text{avg loss}}$; MACD $= \text{EMA}_{12} - \text{EMA}_{26}$; Bollinger Bands $= \mu \pm 2\sigma$ over rolling window.
  \item Airflow DAG design: idempotent tasks, catchup=False for live data, task retries with exponential backoff.
  \item Database schema: time-series table with (symbol, timestamp, OHLCV, indicators); partitioned by date for query performance.
  \item Screening rules: configurable filter chains; e.g.\ RSI $< 30$ AND price $>$ 50-day MA AND volume spike.
\end{itemize}

\section{Data Flow}

Airflow DAG trigger → Yahoo Finance REST → OHLCV normalisation →
Technical indicator calculation → Screening rule application →
PostgreSQL upsert → FastAPI read → React display.


\section{Technology Stack}

\begin{tabular}{ll}
\toprule
\textbf{Component} & \textbf{Technology} \\
\midrule
Orchestration & Apache Airflow 2.x \\
Data sources & yfinance, Alpaca Market Data API \\
Indicators & pandas-ta, TA-Lib \\
Database & PostgreSQL \\
API & FastAPI \\
AI analysis & LiteLLM \\
\bottomrule
\end{tabular}


\section{Design Decisions}

\begin{itemize}
  \item Two separate DAGs (intraday vs.\ nightly) to isolate latency-sensitive
    real-time paths from batch operations.
  \item yfinance as primary source (free); Alpaca as fallback for real-time data
    requiring an API key.
\end{itemize}


\section{Repository Structure}
\begin{lstlisting}[language=bash,caption={Top-level directory layout}]
market_screener/
  .github/
    copilot-instructions.md   # GitHub Copilot guardrails
    agents/
      ecr-secrets-agent.agent.md  # ECR secret rotation runbook
  CLAUDE.md                   # Claude Code guardrails
  GEMINI.md                   # Gemini CLI guardrails
  INFRASTRUCTURE.md           # Operational runbook
  ROADMAP.md                  # Execution log and roadmap
  README.md                   # Project overview
\end{lstlisting}

\section{Contributing}
\begin{enumerate}
  \item Read \texttt{README.md}, \texttt{INFRASTRUCTURE.md}, and \texttt{ROADMAP.md} before any task.
  \item Follow the Code Documentation Standard: every public function must have
    a docstring (Google-style for Python, JSDoc for TypeScript, XML for C\#).
  \item Run the guardrails validator before committing:
    \begin{lstlisting}[language=bash]
git add -A && ./.guardrails/bin/validate_guardrails.sh
    \end{lstlisting}
  \item For deployment or destructive operations, always use a named tmux session:
    \begin{lstlisting}[language=bash]
tmux new-session -A -s deploy-market_screener
    \end{lstlisting}
\end{enumerate}

\end{document}
